\section{Backend Trade-off: Noir+Honk vs.\ Plonky3}
\label{app:backend}

This appendix compares the deployed Noir/UltraHonk recovery circuit
(Section~\ref{subsec:reanchoring}) with a STARK-based alternative, Plonky3.
The comparison is necessarily partial: boundary-signature verification,
condition~(e), dominates the deployed circuit's cost (1{,}097{,}274 opcodes
with it versus 6{,}143 without, a factor of $\sim$180), and no Plonky3
configuration we built implements it.

\begin{table}[h]
\caption{Deployed Circuit: Constant-Cost Scaling (Noir+Honk, Poseidon2)}
\label{tab:reanchor-scaling}
\centering
\footnotesize
\begin{tabular}{@{}rrrr@{}}
\toprule
\textbf{Interval count} & \textbf{Prove time} & \textbf{Verify time} & \textbf{Proof size} \\
\midrule
1   & 38.3 s & 25 ms & 14{,}656 B \\
16  & 41.4 s & 23 ms & 14{,}656 B \\
256 ($N_{\max}$) & 40.4 s & 24 ms & 14{,}656 B \\
\bottomrule
\end{tabular}
\end{table}

Table~\ref{tab:reanchor-scaling} confirms that proving time, verification
time, and proof size are independent of the number of real headers up to
the compiled capacity $N_{\max}{=}256$. Verification by the validators'
embedded WebAssembly verifier costs $\sim$26\,ms per proof. An interval of
$L$ blocks requires $\lceil L/256 \rceil$ leaf proofs when submitted one by one,
or a single aggregate proof up to 1024 blocks (Appendix~\ref{app:aggregation}).

\begin{table}[h]
\caption{Backend Comparison, Statement Coverage}
\label{tab:backend-coverage}
\centering
\footnotesize
\begin{tabular}{@{}lccccc@{}}
\toprule
\textbf{Config.} & (a) & (b) & (c) & (d) & (e) \\
\midrule
Noir+Honk (deployed) & \checkmark & \checkmark & \checkmark & \checkmark & \checkmark \\
Plonky3, primitive-only & -- & -- & -- & -- & -- \\
Plonky3, fuller statement & \checkmark & \checkmark & \checkmark & \checkmark & -- \\
\bottomrule
\end{tabular}
\\[2pt]
\footnotesize (a) chain continuity, (b) FSM/withdrawal-lock policy legality,
(c) circuit-breaker adherence, (d) state-root binding, (e) boundary
ed25519-signature verification.
\end{table}

Table~\ref{tab:backend-coverage} makes the coverage gap explicit. Under a
primitive-only comparison (256 chained Poseidon2 permutations, no FSM or
signature logic, i.e.\ neither configuration implements condition~(e)),
Plonky3 produces a substantially larger proof (1{,}278{,}939~B vs.\
14{,}656~B, $\sim$87$\times$) but proves faster (0.044~s vs.\ 0.684~s,
$\sim$15.5$\times$) with a transparent, post-quantum-secure setup versus the
Noir/Honk backend's trusted KZG setup. Under a fuller-statement comparison
that adds FSM legality and root binding (but still omits condition~(e)),
Plonky3 proves in 0.024--0.096\,s with proof sizes of 84--124\,KB across
2--256 headers. Neither comparison supports a general speed claim for Plonky3, because
both omit the component that dominates the deployed circuit's cost. These
numbers characterize primitive-level backend trade-offs only.

\FloatBarrier
